\chapter{Architecture decisions}\label{ch:decisions}
\chapterstatus{complete}{2026-11-17}{pull requests \#14–\#19 and the analytics change, \src{infra/aws}}

Each record states a choice, what it was chosen over, what it costs, and when to look at it
again. Records are never deleted: a later decision \emph{supersedes} an earlier one and says so.
Add one for every deliberate trade-off (AI-DLC's design stages propose them).

% \adr{number}{title}{date}{status}{context}{decision}{consequences}{revisit when}
\newcommand{\adr}[8]{%
  \section[ADR-#1: #2]{ADR-#1\quad #2}\label{adr:#1}
  {\sffamily\small\color{muted}#3\quad·\quad #4}\par\smallskip
  \begin{description}
    \item[Context] #5
    \item[Decision] #6
    \item[Consequences] #7
    \item[Revisit when] #8
  \end{description}}

\adr{001}{Containers on ECS Fargate, ARM64}{2026-11-16}{Accepted}
  {Slotty is one Next.js app plus a long-running worker that holds a PostgreSQL \code{LISTEN}
   session. It must be cheap at bootstrap, elastic, and free of servers to patch.}
  {Run the existing image on ECS Fargate with Graviton (ARM64) tasks.}
  {No servers to manage; per-second billing; about 20\,\% cheaper than x86. Rejected: EKS (its
   control plane alone costs about \$73 a month), Lambda (cannot hold the worker's session or run
   it continuously), a single EC2 or Lightsail host (not elastic, one more thing to patch).}
  {Steady load makes a reserved EC2 capacity provider cheaper than Fargate.}

\adr{002}{API Gateway HTTP API instead of a load balancer}{2026-11-16}{Accepted}
  {An Application Load Balancer costs about \$18 a month before traffic, which with the rest of the
   stack exceeds the \$50 budget.}
  {CloudFront → API Gateway HTTP API → VPC link → Cloud Map → web tasks.}
  {About \$1–2 a month. Limits: 30\,s per request, 10\,MB responses, responses buffered rather than
   streamed; API Gateway rewrites the forwarded host, which the app corrects (ADR-009). Cloud Map
   adds a private DNS zone (\$0.50).}
  {A page needs streaming or long requests, or the budget allows an ALB (an \code{alb.tf} beside \code{api.tf}; not written yet).}

\adr{003}{NAT instance instead of a NAT gateway}{2026-11-16}{Accepted}
  {Tasks in private subnets need outbound access (SES, push services, ECR). A managed NAT gateway
   costs about \$33 a month plus \$0.045 per GB; interface endpoints cost about \$7 each per AZ.}
  {fck-nat on one \code{t4g.nano} with EC2 automatic recovery; an S3 gateway endpoint (free) for
   image layers.}
  {About \$7 a month. Single point of failure for egress only: while it is down, users are still
   served but email and image pulls wait.}
  {Egress traffic grows, or an egress outage becomes unacceptable (\tfvar{nat_mode = "gateway"}).}

\adr{004}{Single-AZ RDS \texttt{db.t4g.micro}}{2026-11-16}{Accepted}
  {The database is the largest fixed cost. Aurora Serverless v2 would never scale to zero because
   the worker polls, and costs at least \$43 a month.}
  {RDS PostgreSQL 17, single AZ, with 7-day point-in-time recovery and deletion protection.}
  {About \$14 a month. An availability-zone outage takes the database down until it recovers or
   is restored.}
  {Demo-day downtime would cost more than Multi-AZ (\tfvar{db_multi_az}, about \$14 more).}

\adr{005}{A dedicated account with guardrail SCPs}{2026-11-16}{Accepted}
  {The existing AWS account hosts other workloads; Slotty needs its own bill, blast radius and
   hard limits on spending.}
  {An AWS Organization with a \code{slotty-prod} member account, SCPs (one region, small
   instance classes, no expensive services or purchases, no root user), a \$50 budget and anomaly
   alerts held in the management account.}
  {Mistakes and leaked credentials cannot create expensive resources. Growth needs a reviewed
   SCP change. The Organization cannot be deleted while member accounts exist.}
  {A second environment (staging) is needed: add an account under the same OU.}

\adr{006}{SSM Parameter Store for secrets}{2026-11-16}{Accepted}
  {Secrets Manager costs \$0.40 per secret per month and its rotation features are not needed yet.}
  {SecureString parameters under \code{/slotty/}, generated by Terraform, injected by ECS.}
  {Free. Secrets are in the Terraform state, so the state bucket is private and encrypted. No
   automatic rotation.}
  {Automatic database credential rotation becomes a requirement.}

\adr{007}{Direct database sessions, no proxy}{2026-11-16}{Accepted}
  {The worker wakes on PostgreSQL \code{NOTIFY}; a transaction-mode pooler or RDS Proxy breaks
   \code{LISTEN}.}
  {Tasks connect directly, with a pool of at most 10 connections per process.}
  {Four web tasks and the worker stay well under a \code{db.t4g.micro}'s limit.}
  {Web tasks multiply beyond what the database's connection limit allows.}

\adr{008}{SES over SMTP with IAM user credentials}{2026-11-16}{Accepted}
  {The worker already sends with nodemailer over SMTP; SES's SMTP interface needs IAM user
   credentials, while its API could use the task role.}
  {Keep SMTP; a send-only IAM user restricted to \code{From} addresses on the domain.}
  {No code change. Long-lived credentials exist and must be rotated.}
  {Next change to the mailer: switch to the SES API with the task role and delete the user.}

\adr{009}{Origin lock in the application}{2026-11-16}{Accepted}
  {With no WAF and a public API Gateway endpoint, requests could skip CloudFront (its TLS
   settings, headers and caching). HTTP APIs have no resource policies.}
  {CloudFront sends a secret header; \src{src/proxy.ts} refuses requests without it and restores
   the public host for Server Actions.}
  {Costs nothing. The secret lives in SSM and CloudFront's configuration; rotating it means
   updating both, then redeploying.}
  {A WAF or an ALB with a CloudFront-only security group makes it redundant.}

\adr{010}{CI owns the running image}{2026-11-16}{Accepted}
  {Terraform and CI both want to define task definitions.}
  {Services ignore task-definition changes in Terraform; CI registers revisions from the latest one
   with only the image changed.}
  {CI's role stays narrow. A configuration change in Terraform goes live on the next deploy.}
  {—}

\adr{011}{Immutable OIDC subject for deploys}{2026-11-16}{Accepted}
  {The repository issues GitHub's immutable OIDC subjects, which include owner and repository IDs.}
  {Trust exactly \code{<prefix>:ref:refs/heads/master} with the ID-based prefix.}
  {A renamed or re-created repository with the same name cannot deploy.}
  {The repository moves to another owner.}

\adr{012}{A separate image for migrations}{2026-11-16}{Accepted}
  {The Prisma CLI adds about 270 MB; web tasks pull their image on every start.}
  {A \code{migrate} build target with the CLI and OpenSSL; the app image is 448 MB instead of
   1.97 GB. The migrate image trusts only the RDS CA bundle (\envvar{SSL_CERT_FILE}), because the
   migration engine's \code{sslcert} option reads only the first certificate of a bundle.}
  {Two images per deploy, both tagged with the commit.}
  {Migrations move to a different tool.}

\adr{013}{Fargate Spot for the worker and extra web tasks}{2026-11-16}{Accepted}
  {Spot is about 70\,\% cheaper. The worker tolerates interruption; one web task on regular
   Fargate keeps the site up if Spot capacity disappears.}
  {Web: first task Fargate, the rest Spot. Worker: Spot.}
  {If Spot capacity is unavailable, emails wait and the web tier runs on one task.}
  {Spot interruptions are seen on demo days.}

\adr{014}{Documentation in LaTeX; changes through AI-DLC}{2026-11-16}{Accepted}
  {Experts must be able to audit and extend the system; Markdown files had grown into tables that
   were hard to keep coherent, and changes had no recorded reasoning.}
  {One LaTeX document with a shared design and diagrams, built in a pinned TeX Live image and in
   CI; AWS's AI-DLC for changes, with gates that require the affected chapters to be updated.}
  {A heavier toolchain (a 3 GB build image, used only for the documentation). The old Markdown
   files were folded into \cref{ch:journeys,ch:rules,ch:roadmap}.}
  {The document stops being kept up to date: then prefer fewer chapters over stale ones.}

\adr{015}{First-party usage analytics and anonymous feedback}{2026-11-17}{Accepted}
  {Before launch to about 300 students and 15 staff, improvements should rest on evidence: where
   booking and marking get stuck, how fast pages are on phones, what breaks, what nobody uses, and
   what people say. The budget has about \$2 a month left; the users are students whose data
   includes marks.}
  {Collect it ourselves: a small browser tracker beaconing a fixed event catalog to
   \route{/api/events}, stored in PostgreSQL under a pseudonymous HMAC visitor key with route
   templates only, rolled up daily by the worker, kept 90 days raw and 13 months as totals, and
   reported to the product team (ADR-016). Workflow numbers (reschedules, no-shows, time to mark) are
   queried from the existing tables. Feedback and one-question surveys are stored without the
   user. Usage data is on by default with an opt-out in Account, explained at \route{/privacy};
   errors and feedback are collected either way. A report-only content security policy limits
   the page to its own origin.}
  {About \$0.05 a month; no new AWS service or third party, and the data can be joined with
   Slotty's own (role, flow, release). In exchange we build and maintain the reports ourselves,
   there is no session replay or heat map, and blocked or closed tabs can lose their last few
   events. Rejected: CloudWatch RUM (about \$0.50 a month at this scale, but it needs a Cognito
   identity pool and its data can't be joined with ours), a SaaS such as PostHog or Plausible
   (student data leaves our account; consent banner; vendor lock-in), self-hosted PostHog (needs
   more memory than the whole budget buys), session replay (records marks and personal data).}
  {Traffic grows past about 5 million events a month (move raw events out of RDS), the team
   needs replay or heat maps, or the policy for student data changes (for example opt-in
   consent).}

\adr{016}{Usage insights belong to the product team, not to administration}{2026-11-17}{Accepted}
  {Admins run the deployment (users, courses, email). Usage data and feedback are evidence for
   whoever builds Slotty. Putting them in the admin console widens who can see them, and any
   bug in admin or course code would sit one step from internal data.}
  {A separate \route{/platform} area and permission, \code{User.platformAccess}, checked by
   \code{assertPlatform}. A platform account is dedicated: never an admin and never a course
   member, enforced by \code{canGrantPlatform} and every enrolment path. It is granted only by a
   command-line job (the same way the first admin is created), and the admin console neither
   lists nor changes it, nor shows its activity in the audit log.}
  {Two kinds of console account to operate; the product team needs a second login if they also
   teach or administer. A university admin can't see adoption numbers unless the product team
   shares them. Rejected: an extra flag for admins (keeps the data one bug away from the
   console), a separate dashboard tool such as Metabase or Grafana (its own compute, over the
   budget, and a network path into the private database), a separate deployment (doubles the
   cost).}
  {Admins need adoption numbers for their own courses: add an aggregate, course-level view for
   them rather than opening \route{/platform}.}

